\section{Metodologia}

A presente etapa do projeto concentrou-se na consolidação teórica e na
implementação de uma prova de conceito computacional voltada à aplicação
de algoritmos quânticos variacionais na simulação de sistemas moleculares.
Dessa forma, a metodologia adotada foi estruturada de modo a integrar
capacitação conceitual, domínio de ferramentas computacionais e
implementação prática do algoritmo Variational Quantum Eigensolver (VQE),
permitindo a compreensão completa do fluxo híbrido quântico-clássico
característico dos algoritmos da era NISQ.

\subsection{Consolidação teórica e capacitação técnica}

Inicialmente, realizou-se um processo sistemático de estudo dos fundamentos
da computação quântica, algoritmos variacionais e métodos de simulação
molecular. Essa etapa constituiu parte central das atividades desenvolvidas,
uma vez que o objetivo principal desta fase do projeto consistiu na
consolidação conceitual necessária para futuras aplicações em problemas
de interesse na área de materiais energéticos.

Esta etapa foi realizada por meio 
da participação na disciplina ``Computação Quântica'', ofertada na Unidade Acadêmica de Belo 
Jardim (UFRPE) e ministrada pelo orientador do projeto, Prof. Dr. Fábio 
Magalhães de Novaes Santos. 
Na disciplia, estudamos os fundamentos matemáticos da mecânica quântica, incluindo 
notação de Dirac e representação vetorial de estados 
quânticos, até algoritmos quânticos fundamentais, como os algoritmos de 
Grover e Shor, além de tópicos avançados relacionados à simulação quântica e 
caminhadas quânticas, por exemplo. Ao final da disciplina, apresentou-se um seminário sobre VQE para toda a turma, com a entrega de códigos Qiskit que implementam a versão padrão desse algoritmo para a molécula de H$_2$.
% Essa formação ocorreu simultaneamente ao 
% desenvolvimento da iniciação científica, contribuindo diretamente para a 
% compreensão dos modelos computacionais empregados neste trabalho.

Na disciplina, foram utilizados como referências principais textos clássicos da área de
computação quântica \cite{nielsen2000quantum,hidary2019quantum}, bem como
artigos de revisão sobre algoritmos variacionais e computação quântica
aplicada à química molecular \cite{tilly2022variational,cao2019quantum}. Paralelamente, foi realizado o acompanhamento
do curso educacional \textit{Quantum Chemistry with VQE}, disponibilizado
pela plataforma IBM Quantum Learning \cite{ibm_vqe_course}, o qual apresenta
de forma estruturada as etapas necessárias para implementação do algoritmo
VQE em problemas de estrutura eletrônica. 
Essa etapa permitiu o domínio conceitual dos elementos fundamentais do
pipeline variacional, incluindo construção do Hamiltoniano molecular,
definição de estados variacionais, estratégias de medição e métodos de
otimização clássica.

\subsection{Ambiente computacional e ferramentas utilizadas}

Conforme previsto no plano de trabalho inicial, foram consideradas as 
plataformas \textcite{quantumket}, QuBox e Qiskit para a realização das simulações quânticas. 
Entretanto, durante a fase de implementação, observou-se que o ecossistema 
desenvolvido pela IBM Quantum, baseado no framework Qiskit, apresenta 
integração completa entre simulação local, execução em hardware quântico 
real e ferramentas de química quântica variacional. Pretende-se posteriomente comparar esse framework com a plataforma nacional composta de Ket e QuBox.

Dessa forma, optou-se pela adoção prioritária do Qiskit como ambiente 
unificado de desenvolvimento, permitindo a execução de todas as etapas 
necessárias ao projeto sem perda metodológica ou experimental. Essa escolha 
reduziu a complexidade do fluxo computacional e favoreceu uma maior 
padronização dos experimentos realizados.

As simulações foram executadas exclusivamente em simuladores quânticos
clássicos disponibilizados pelo Qiskit Aer, possibilitando a análise
do comportamento ideal do algoritmo na ausência de ruídos experimentais.
Essa escolha permitiu isolar aspectos conceituais e metodológicos do
algoritmo, priorizando a compreensão do seu funcionamento interno. Pretende-se realizar testes em hardware quântico real nas fases posteriores do projeto.

\subsection{Formulação do problema molecular}

Foram selecionadas como sistemas modelo as moléculas de hidrogênio (H$_2$) e 
hidreto de lítio (LiH), amplamente utilizadas na literatura como benchmarks 
em simulações quânticas devido à sua simplicidade estrutural e relevância na 
validação de algoritmos variacionais. A metodologia utilizada segue o tutorial do curso de VQE da \textcite{ibm_vqe_course}.

\subsubsection{Definição geométrica e base atômica}

As geometrias moleculares foram definidas por meio de coordenadas 
cartesianas dos núcleos atômicos. Para a molécula de hidrogênio, foi adotada 
uma distância internuclear de 0,735 Å, valor próximo ao comprimento de 
ligação de equilíbrio experimental. 
Em ambas as moléculas foi utilizada a base atômica mínima STO-3G, escolha 
amplamente empregada em estudos introdutórios por reduzir significativamente 
o custo computacional mantendo as principais características físicas do 
sistema eletrônico.

\subsubsection{Construção do Hamiltoniano eletrônico}

A construção do Hamiltoniano eletrônico em segunda quantização foi realizada 
por meio do módulo \textit{Qiskit Nature}, utilizando o \textit{PySCFDriver} 
como interface para cálculo da estrutura eletrônica inicial via Hartree–Fock 
restrito (RHF).
O operador fermiónico resultante foi automaticamente gerado pelo framework e 
posteriormente mapeado para operadores de spin por meio da transformação de 
Jordan–Wigner, permitindo sua representação como combinação linear de 
operadores de Pauli mensuráveis em dispositivos quânticos digitais.
Para fins de comparação e validação dos resultados variacionais, a energia 
exata do estado fundamental foi obtida utilizando o algoritmo clássico 
\textit{NumPyMinimumEigensolver}, equivalente ao cálculo de FCI no espaço 
eletrônico considerado.

\subsection{Implementação do algoritmo Variational Quantum Eigensolver}

\subsubsection{Definição do \textit{ansatz} variacional}

Foi adotado um \textit{ansatz} do tipo \textit{hardware-efficient}, 
implementado por 
meio da classe \textit{EfficientSU2} do Qiskit. Esse tipo de circuito 
consiste em camadas de rotações parametrizadas intercaladas com portas de 
entrelaçamento total (\textit{full entanglement}), sendo adequado ao regime 
NISQ por apresentar profundidade moderada e boa capacidade expressiva.
O número de repetições do bloco variacional foi fixado em duas camadas 
(\textit{reps}=2), buscando equilíbrio entre expressividade do circuito e 
estabilidade do processo de otimização.

\subsubsection{Loop híbrido quântico-clássico}

O algoritmo VQE foi implementado utilizando a arquitetura híbrida padrão:
\begin{itemize}
\item preparação do estado variacional parametrizado;
\item estimativa do valor esperado da energia por meio de simulação ideal de 
estado vetorial;
\item atualização clássica dos parâmetros variacionais;
\item repetição do processo até o critério de parada, após 200 iterações.
\end{itemize}
A estimativa do valor esperado foi realizada utilizando o 
\textit{StatevectorEstimator}, garantindo cálculo determinístico da energia 
sem ruído estatístico.

\subsubsection{Otimização clássica}

A minimização da energia foi conduzida por meio do otimizador clássico 
COBYLA, configurado com limite máximo de 200 iterações. O critério de parada 
adotado baseou-se no número máximo de iterações definido, sendo 
adicionalmente monitorada a estabilização da energia ao longo do processo 
iterativo.

\subsection{Execução das simulações e análise dos dados}

As simulações foram inicialmente realizadas para a molécula de hidrogênio 
(H$_2$), com o objetivo de validar a implementação do fluxo híbrido. 
Posteriormente, o procedimento foi estendido ao hidreto de lítio (LiH), 
permitindo avaliar o comportamento do algoritmo diante de um sistema 
molecular de maior complexidade eletrônica.
Durante a execução do VQE, os valores da energia a cada iteração foram 
registrados por meio de função de \textit{callback}, possibilitando a 
análise detalhada da convergência do algoritmo.

Os resultados finais, incluindo energia variacional, energia exata de 
referência (FCI), erro absoluto e número de iterações, foram armazenados em 
arquivos no formato CSV, garantindo reprodutibilidade e organização 
sistemática dos dados para posterior geração de gráficos de convergência.
A análise concentrou-se na taxa de convergência do algoritmo, estabilidade 
do otimizador e precisão da energia estimada em relação ao valor exato, 
caracterizando o estudo como uma prova de conceito do fluxo computacional do 
VQE em ambiente controlado e ideal.